% ch3 D.1 前段 激子 (原书页 132–138 / PDF p147–p153)
%--C-TAIL--
%JOIN%
思维的种种基础。尽管如此，人们也可能过分强调处理固体时所能达到的、乃至宜于讲求的那种严格性。正如 Landau 在旧版《Statistical Physics》的序言中所说：“我们在这里谈的是理论物理，因此数学上的严格性当然是既不相干、也办不到的。”这话并不完全正确，但也相去不远。

我们最好记住：第一，我们用来工作的数学模型，几乎总是对实际物理系统已经作了高度简化的理想化；第二，准粒子物理就其本性而言全部建立在微扰级数之上，因为它无非是试图用一个形如仅有弱相互作用粒子系统的理论去逼近一个完整的理论。但我们必须时时猜测微扰论应当从什么状态出发——也就是说，猜测系统的物理本性究竟是什么：金属型的、磁性的、绝缘的、液态的还是固态的。这种猜测可能正确，但也可能错得离谱；在后一种情形，往往唯一的迹象只是我们的微扰论出现轻微的、几乎觉察不到的发散。超导电性里发生的正是这种事。第三点警告是：由于存在集体激发，几乎所有微扰级数都确实会在至少某个方面失效，而这些集体激发又往往只能作为各项发散子级数之和从微扰论中求得。总的说来，我认为下述态度至少是站得住脚的：固体物理的奠基者们和 Landau 学派对准粒子的本能式使用，其严格性比起精巧的多体理论几乎毫不逊色，而且也许更万无一失，因为它的\emph{外表}不那么严格。

\section{集体激发}

\subsection{激子}

绝缘体中的激子（exciton）与铁磁体中的自旋波（spin wave）是最简单的两类集体激发（collective excitation）。它们在许多方面惊人地相似，在与它们相联系的形式数学方面尤其如此；两者还有更深一层的相似之处：它们的零点运动（zero-point motion）都相对很少，由零点运动所带来的复杂问题也很少。在每一种情形下，原因都在于：正如在 $N+1$ 体问题中那样，人们可以从一个几乎能够严格定义的基态出发来表述问题。在铁磁体中之所以如此，是因为铁磁基态的特征（如果它存在的话）是自旋的总 $Z$ 分量为 $M_Z = N/2$。所有自旋波态虽然能量可能非常低，其特征却是 $M_Z = (N/2) - 1$，因而把它们与基态联系起来的哈密顿量矩阵元充其量也非常小。

激子的情形则更为复杂。如读者所知，激子是光激发，我们既可以按照 Frenkel (77) 的模型、也可以按照 Wannier (78) 的模型来讨论它们。Frenkel 模型适用于紧束缚的情形，例如分子型或离子型晶体，它由一个在晶体中移动的受激原子态构成。例如，它可以是这样一种 NaCl 中的态：Na$^+$ 离子实的 p 电子之一被激发到一个 s 能级，而由此形成的激发可以在晶体中移动。Wannier 模型更常被观察到，或许也更有趣；在这一模型中，人们假定一个电子从一个能带被激发到下一个能带，但由此形成的准粒子与准空穴按弱势 $e^2/\kappa R$ 相互吸引，并一起在晶体中穿行。无论哪种情形，通常都相当清楚该如何定义一个不包含任何激子的固定基态；例如，对光学活性的（optically active）激子——诸如 $sp$ 跃迁——$k=0$ 的激子态与基态宇称相反，因而不能与基态混合。然而我们将会看到，这一简化并非对所有激子都完全成立。

我这里的材料将主要取自 J.~J.~Hopfield (79) 的一篇出色的论文，以及 E.~Dresselhaus (80) 的另一篇。我不认为存在任何好的综述；较早的理论在 Seitz 的书中有论述。

就定性物理图象的目的以及与自旋波作比较而言，Frenkel 模型更为有用。让我们从它出发，然后再从它出发作微扰。

Frenkel 模型（至少在一开始）所设想的是一种像固态 Ne 或 Ar 那样的晶体，其中的原子相距相当远，几乎互不交叠。当然，这种晶体是一种绝缘体。我们暂且撇开对基本多体问题的任何讨论，而用 Hartree-Fock 方法处理势场；相关的相互作用我们稍后将用多体理论来处理，但从 Hartree 出发最为简单，同时按惯例附加上一个条件：由于我们处理的是元激发，多体效应并不太严重。最后一个被填满的能带，我们假定是一个 s 带。把这个带的 Wannier 函数记作 $a_s(r-R_j)$。由于这个带极窄，非对角矩阵元
\begin{equation*}
b_{jk}^{s} \;=\; \int a_s^{*}(R_j - r)\; \mathcal{H}_{1\mathrm{el}}(r)\; a_s(r - R_k)\; dr
\end{equation*}
非常小——至于相对于什么而言小，我们很快就会看到。

还有一个“导”带——一个由激发态构成的能带，为简单起见我们同样假定它具有 p 型 Wannier 函数
\begin{equation*}
a_p^{\,x,y,z}(r - R_j)
\end{equation*}
以及非常窄的带宽（即转移积分，transfer integrals）
\begin{equation*}
b_{jk}^{xx'},\ \text{等等}
\end{equation*}
自由原子会有一定的激发能
\begin{equation*}
E_0 \;=\; (E_p - E_s)
\end{equation*}
——从 s 态激发到 p 态——在这样一种固体中，它不会由于变换到 Wannier 函数而受到太严重的修正。必须认识到，这个激发能在数量上并不接近 s 带与 p 带之间的 Hartree-Fock 平均能量差。

其原因在于：能带能量差对应的态，是电子被完全激发而离开自己原来那个原子、从而发现自己处在一个既有 s 又有 p 电子的原子上的态，因此除了单电子激发能 $E_0$（图 42）之外，还有如下的额外库仑排斥能
\begin{equation*}
U \;=\; \int |a_p(r-R_j)|^2\; \frac{e^2}{|r-r'|}\; |a_s(r'-R_j)|^2\; dr\; dr'
\end{equation*}

%FIG{42}: 图 42（能级示意图：上方为 E_k(p) 能带曲线，下方为 E_p(atomic) 原子能级线组与 E_S 基态能级线；竖直箭头分别标出 E_0、E_k(p)−E_k(s) 两个激发能以及 U 的间隔）

换一种说法：电子与它身后留下的空穴通过完全的库仑相互作用（仅被介质的介电常数稍微修正）彼此强烈吸引，因此电子被激发到同一原子上的 p 态，远比把它激发到邻近的原子上去容易得多。

现在很清楚，如果各个 $b$ 与 $U$ 相比很小，那么 Frenkel 理论将相当好地成立，因为无论空穴还是电子都不会倾向于跳离自己的同伴。这时激子
%FIG{exciton-jkj}: 无编号正文插图（原书不编号）——Frenkel 激子组态示意：上排为 p 能级，格点 j 处有一个电子（⊖），一道弧形箭头从该电子指向右侧格点 k 处的空 p 能级，箭头旁标 “+ U”；下排为 s 能级，格点 j 处为空穴，k 与 j' 处各有一个电子（⊖）
将是紧束缚的，而且在 s 空穴-p 电子带态的连续谱与激子之间将有一个很大的能隙。在另一种情形 $U < b$ 中，电子与空穴将比较容易分开，Wannier 图象会好得多。我们将在后面讨论那种情形。

另一方面，即使电子与空穴耦合得很紧，这一对粒子仍然能够一起在晶体中移动。在我们现在的情形中，实现这一点最重要的机制，是电子与相邻原子上的电子之间的相互作用。这个相互作用可以写成
\begin{equation*}
\int \psi^{\dagger}(r)\,\psi^{\dagger}(r')\,\frac{e^2}{|r-r'|}\,\psi(r')\,\psi(r)\,dr\,dr' \tag{11}
\end{equation*}
其中 $\psi$ 可以用 Wannier 函数的产生与湮灭算符展开，即 $c_j^s$ 与 $c_j^{p_x}$、$c_j^{p_y}$、$c_j^{p_z}$：
\begin{equation*}
\psi^{\dagger}(r) \;=\; \sum_j \sum_{n=s,\,p_x,\,p_y,\,p_z} c_j^{n\dagger}\, a_n^{*}(r-R_j)
\end{equation*}
有各种各样的项可能产生把一个原子激发、同时使其邻居退激发的效果；其中最重要的是“直接的”（direct）即偶极（dipolar）项，其典型的一员是
\begin{equation*}
\left(c_j^{p_x}\right)^{\dagger} c_k^{s\dagger} c_k^{p_x} c_j^{s} \int a_{p_x}^{*}(r-R_j)\, a_s(r-R_j)\, \frac{e^2}{|r-r'|}\, a_s^{*}(r'-R_k)\, a_{p_x}(r'-R_k)\, dr\, dr'
\end{equation*}
这类积分用多极展开（multipolar expansion）来求值最为容易：
\begin{align*}
\frac{e^2}{|r-r'|} = \; & e^2 \Bigg[ \frac{1}{R_j - R_k} + \left\{ (r-R_j) + (r'-R_k) \right\} \cdot \nabla\!\left(\frac{1}{r}\right)\Big|_{R_j-R_k} \\
& + (r-R_j)(r-R_k) \cdot \nabla\nabla\!\left(\frac{1}{r}\right)\Big|_{R_j-R_k} + \cdots \Bigg] \tag{12}
\end{align*}

（译注：式 (12) 照录原书。按标准多极展开，第二项花括号内应为 $\{(r-R_j)-(r'-R_k)\}$，即 $(r'-R_k)$ 之前应为负号；原书印作正号。）

第一个不为零的项是偶极项，即上式最后一项，它给出的典型矩阵元为
\begin{equation*}
|X_{sp}|^2 \cdot T(r_j - R_k)_{xx}
\end{equation*}
其中 $X_{sp}$ 是原子态的偶极矩矩阵元：
\begin{equation*}
X_{sp} \;=\; e \int a_{p_x}^{*}(r)\; x\; a_s(r)\; dr
\end{equation*}
而 $T$ 是偶极相互作用算符
\begin{equation*}
T_{jk} \;=\; \frac{1}{R_{jk}^{3}} \left\{ 1 - 3\,\frac{\mathbf{R}_{jk}\,\mathbf{R}_{jk}}{R_{jk}^{2}} \right\}
\end{equation*}

在这个偶极相互作用算符——它表示与每个原子相联系的运动电偶极子之间的直接静电相互作用——之中，我上面给出的那种形式的矩阵元所联系的各个态，正如我先前概略说明的那样，每一个都比基态高出大约 $E_0$ 的能量。还有一些形如 $c_j^{p\dagger} c_j^{s} c_k^{p\dagger} c_k^{s}$ 的项，它们同时产生或湮灭两个激发；这些项有相当的重要性，我们很快就会回头来讨论它们，但就目前而言，它们显然不如其他项有效，因为它们的能量分母 $\Delta E$ 为 $2E_0$。于是我们的哈密顿量中现在有
\begin{align*}
\mathcal{H}_{\mathrm{unpert}} = \sum_{j,\,p_i} \frac{E_0}{2} & \left[ \left(c_j^{p_i}\right)^{\dagger} c_j^{p_i} - c_j^{s*} c_j^{s} \right] \\
& + \frac{1}{2} \sum_{j \neq k}\; \sum_{p_i\, p_{i'}} \Bigl( \left(c_j^{p_i}\right)^{\dagger} c_k^{s\dagger} c_k^{p_{i'}}\, c_j^{s}\; |X_{sp}|^2\; T_{jk} \Bigr)_{ii'} \tag{13}
\end{align*}

现在有用的一步，是几乎在集体激发理论的每一个问题中都要采取的一步：从费米子（fermion）算符对到玻色子（boson）的变换。这类变换中人们最熟悉的，是电子自旋算符这一特殊情形。假定代替三个 p 函数的只是唯一的一个 $\phi_p$。那么出现在哈密顿量——式 (13)——之中、与原子 $j$ 相联系的算符，将是如下三种组合
\begin{align*}
c_{j1}^{\dagger}\, c_{j1} - c_{j0}^{\dagger}\, c_{j0} \qquad & \qquad 0 = s \\
c_{j1}^{\dagger}\, c_{j0} \;\; \text{与} \;\; c_{j0}^{\dagger}\, c_{j1} \qquad & \qquad 1 = p
\end{align*}
此外，我们只对原子 $j$ 要么被激发、要么未被激发的那些态感兴趣——也就是说，对 $n_p + n_s = 1$ 的态感兴趣。

这些算符对恰好等价于一组 Pauli 自旋算符
\begin{equation*}
\sigma_{jz} \;=\; \begin{pmatrix} 1 & 0 \\ 0 & -1 \end{pmatrix}\;\begin{matrix} \mathrm{p}_{occ} \\ \mathrm{s}_{occ} \end{matrix} \qquad\quad \sigma_j^{+} \;=\; \begin{pmatrix} 0 & 1 \\ 0 & 0 \end{pmatrix} \qquad\quad \sigma_j^{-} \;=\; \begin{pmatrix} 0 & 0 \\ 1 & 0 \end{pmatrix}
\end{equation*}
——它们作用在描写第 $j$ 个原子状态的赝自旋（pseudo-spin）波函数上。

首先，属于不同原子的费米子算符对彼此对易，正如自旋算符那样；其次，在涉及单个原子的波函数的子空间内，这些自旋算符恰好具有与
\begin{equation*}
c_1^{\dagger} c_1 - c_0^{\dagger} c_0 \;(\sigma_z) \qquad\qquad c_1^{\dagger} c_0 \;(\sigma^{+}); \quad c_0^{\dagger} c_1 \;(\sigma^{-})
\end{equation*}
完全相同的代数性质。因此在那种情形下，把哈密顿量写成
\begin{equation*}
\mathcal{H} \;=\; \frac{E_0}{2} \sum_j \sigma_{jz} \;+\; \frac{1}{4} \sum_{jk} X_j\, T_{jk}\, X_k \left( \sigma_j^{+} \sigma_k^{-} + \sigma_j^{-} \sigma_k^{+} \right)
\end{equation*}
的形式将是精确无误的。在磁学情形中恰好就要用到这种对应。

从自旋算符到玻色子的变换，Dyson (81) 已经作过详尽的讨论；借助一个小技巧可以把它简化，这个技巧的出处我已经记不得了。只要两个自旋算符——或者两个费米子算符对——涉及的是不同的原子，它们就彼此对易，正如玻色子那样。但当两个自旋算符涉及同一个原子时，它们的对易关系要更接近费米型得多，而这正是建立相应的一组玻色子时的主要困难。

在激子的情形以及自旋的许多情形中，这并不是一个很严重的限制，因为出现的激发总数相当少，从而它们相遇重合的可能性可以忽略。在这样的情形下，我们或许可以猜想：按照下述显而易见的对应规则直接把旋量（spinor）认同为玻色子，就会令人满意：
\begin{equation*}
\sigma_j^{\dagger} \rightarrow b_j^{\dagger} \qquad \sigma_j \rightarrow b_j \qquad \frac{\sigma_{jz} + 1}{2} \rightarrow n_j \;=\; b_j^{\dagger} b_j \tag{14}
\end{equation*}
就最重要的那一对能级 $n_j = 0$ 与 $1$ 而言，算符组 (14) 具有完全相同的代数性质。但玻色子 $b$ 却具有与一整串额外能级 $n_j = 2,\, 3,\, \ldots$ 相联系的矩阵元。

使对应 (14) 成为精确的诀窍很简单：切断能级阶梯最底下两级与其余各级之间的联络。插入一道“势垒”（barrier），它恰好把连接 $n_j = 1$ 与 $n_j = 2$ 的矩阵元投影掉（project out），而在其他方面并不改变 $n_j = 0$ 和 $1$ 的矩阵元（图 43）。这道“势垒”就是：在所有使用 $b_j^{\dagger}$ 的地方都在其之前、在使用 $b_j$ 的地方都在其 \emph{之后} 乘以 $(1-n_j)$。也就是说，
\begin{equation*}
\sigma_j^{\dagger} = b_j^{\dagger}(1-n_j) \qquad \sigma_j = (1-n_j)\, b_j \qquad \sigma_{jz} = 2n_j - 1
\end{equation*}

%FIG{43}: 图 43（示意图：左边的 σ_j 只有两个能级（0 与 1）；右边的 b_j 具有包括 n_j = 0、1 与 2、3、4、……在内的整个能级阶梯，在 n_j = 1 与 2 之间画有锯齿状 “Barrier”（势垒））

立刻可以验证：当 $n_j = 1$ 时，$\sigma_j^{\dagger}$ 无法把它升到 $n_j = 2$，反之亦然。

别的技巧也可能同样好，例如 Holstein-Primakoff 展开或 Wentzel 的有效哈密顿量方法 (82)；某些格林函数（Green's function）方案则更为一般 (93)；但就费米子对到玻色子的变换而言，这是迄今最简单的一个论证。利用这一变换，在单激发能级的情形下，哈密顿量现在——仍然是严格地——为
\begin{equation*}
\mathcal{H} \;=\; E_0 \sum_j n_j \;+\; \sum_{jk} b_j^{\dagger}(1-n_j)\, T_{jk}\, (1-n_k)\, b_k\; \frac{|X|^2}{2}
\end{equation*}
把它推广到激发态简并的情形是平凡的。在我们正在讨论的情形中，可以把这对能级——其中之一是三重简并的——换成三维谐振子最低的两个能级
\begin{equation*}
b_j^{\dagger} \;=\; b_j^{x\dagger},\; b_j^{y\dagger},\; b_j^{z\dagger}
\end{equation*}
于是，为了施加我们的势垒，只需把 $b^x$ 乘以 $(1-n_j)$，其中 $n_j = n_{jx} + n_{jy} + n_{jz}$；这样便有 $c_j^{p_x}{}^{\dagger} c_j^{s} = b_j^{x}{}^{\dagger}(1-n_j)$，等等。所得的用玻色子表示的哈密顿量为

% （本文件到此为止；下接原书 p154，由下一文件续译）
